\chapter{2019年全国II卷（文）}

\section{单选题}

\begin{problem}
已知集合 \(A=\{x\mid x>-1\}\)，\(B=\{x\mid x<2\}\)，则 \(A\cap B=\)（\quad）

\choices
    {\((-1,+\infty)\)}
    {\((-\infty,2)\)}
    {\((-1,2)\)}
    {\(\varnothing\)}
\end{problem}

\begin{answer}
C
\end{answer}

\begin{solution}
由 \(x>-1\) 且 \(x<2\)，得 \(-1<x<2\)，所以 \(A\cap B=(-1,2)\)，故选 C.
\end{solution}

\begin{problem}
设 \(z=\i(2+\i)\)，则 \(\overline z=\)（\quad）

\choices
    {\(1+2\i\)}
    {\(-1+2\i\)}
    {\(1-2\i\)}
    {\(-1-2\i\)}
\end{problem}

\begin{answer}
D
\end{answer}

\begin{solution}
由 \(\i^2=-1\)，得 \(z=\i(2+\i)=2\i+\i^2=-1+2\i\)，所以 \(\overline z=-1-2\i\)，故选 D.
\end{solution}

\begin{problem}
已知向量 \(\bs a=(2,3)\)，\(\bs b=(3,2)\)，则 \(|\bs a-\bs b|=\)（\quad）

\choices
    {\(\sqrt2\)}
    {\(2\)}
    {\(5\sqrt2\)}
    {\(50\)}
\end{problem}

\begin{answer}
A
\end{answer}

\begin{solution}
由 \(\bs a-\bs b=(2-3,3-2)=(-1,1)\)，得 \(|\bs a-\bs b|=\sqrt{(-1)^2+1^2}=\sqrt2\)，故选 A.
\end{solution}

\begin{problem}
生物实验室有 \(5\) 只兔子，其中只有 \(3\) 只测量过某项指标. 若从这 \(5\) 只兔子中随机取出 \(3\) 只，则恰有 \(2\) 只测量过该指标的概率为（\quad）

\choices
    {\(\dfrac23\)}
    {\(\dfrac35\)}
    {\(\dfrac25\)}
    {\(\dfrac15\)}
\end{problem}

\begin{answer}
B
\end{answer}

\begin{solution}
从 \(3\) 只测量过该指标的兔子中取 \(2\) 只，从另外 \(2\) 只中取 \(1\) 只的取法数为 \(\mathrm C_3^2\mathrm C_2^1\)，从 \(5\) 只中取 \(3\) 只的取法数为 \(\mathrm C_5^3\). 因此所求概率为 \(\dfrac{\mathrm C_3^2\mathrm C_2^1}{\mathrm C_5^3}=\dfrac{3\times2}{10}=\dfrac35\)，故选 B.
\end{solution}

\begin{problem}
在“一带一路”知识测验后，甲，乙，丙三人对成绩进行预测.

甲：我的成绩比乙高.

乙：丙的成绩比我和甲的都高.

丙：我的成绩比乙高.

成绩公布后，三人成绩互不相同且只有一个人预测正确，那么三人按成绩由高到低的次序为（\quad）

\choices
    {甲，乙，丙}
    {乙，甲，丙}
    {丙，乙，甲}
    {甲，丙，乙}
\end{problem}

\begin{answer}
A
\end{answer}

\begin{solution}
若乙的预测正确，则丙的成绩高于乙和甲，因而丙的预测“我的成绩比乙高”也正确，与只有一个人预测正确矛盾，所以乙的预测错误. 因此甲、丙两人的预测中恰有一个正确：若甲预测正确而丙预测错误，则成绩顺序为甲、乙、丙；若甲预测错误而丙预测正确，则丙的成绩高于乙，从而乙的预测也正确，仍然矛盾. 所以三人按成绩由高到低的次序为甲、乙、丙，故选 A.
\end{solution}

\begin{problem}
设 \(f(x)\) 为奇函数，且当 \(x\ge0\) 时，\(f(x)=\e^x-1\)，则当 \(x<0\) 时，\(f(x)=\)（\quad）

\choices
    {\(\e^{-x}-1\)}
    {\(\e^{-x}+1\)}
    {\(-\e^{-x}-1\)}
    {\(-\e^{-x}+1\)}
\end{problem}

\begin{answer}
D
\end{answer}

\begin{solution}
当 \(x<0\) 时，\(-x>0\)，由奇函数性质得 \(f(x)=-f(-x)\). 因为 \(f(-x)=\e^{-x}-1\)，所以 \(f(x)=-\e^{-x}+1\)，故选 D.
\end{solution}

\begin{problem}
设 \(\alpha\)，\(\beta\) 为两个平面，则 \(\alpha\parallel\beta\) 的充要条件是（\quad）

\choices
    {\(\alpha\) 内有无数条直线与 \(\beta\) 平行}
    {\(\alpha\) 内有两条相交直线与 \(\beta\) 平行}
    {\(\alpha\)，\(\beta\) 平行于同一条直线}
    {\(\alpha\)，\(\beta\) 垂直于同一平面}
\end{problem}

\begin{answer}
B
\end{answer}

\begin{solution}
若 \(\alpha\parallel\beta\)，则 \(\alpha\) 内任意两条相交直线都与 \(\beta\) 平行，故选项 B 的条件成立. 反之，若 \(\alpha\) 内有两条相交直线都与 \(\beta\) 平行，则这两条相交直线的方向确定的平面与 \(\beta\) 平行；该平面就是 \(\alpha\)，所以 \(\alpha\parallel\beta\). 故选 B.
\end{solution}

\begin{problem}
若 \(x_1=\dfrac{\pi}{4}\)，\(x_2=\dfrac{3\pi}{4}\) 是函数 \(f(x)=\sin\omega x\)（\(\omega>0\)）两个相邻的极值点，则 \(\omega=\)（\quad）

\choices
    {\(2\)}
    {\(\dfrac32\)}
    {\(1\)}
    {\(\dfrac12\)}
\end{problem}

\begin{answer}
A
\end{answer}

\begin{solution}
函数 \(f(x)=\sin\omega x\) 的相邻极值点对应的相位相差 \(\pi\). 因而 \(\omega(x_2-x_1)=\pi\)，即 \(\omega\left(\dfrac{3\pi}{4}-\dfrac{\pi}{4}\right)=\pi\)，解得 \(\omega=2\)，故选 A.
\end{solution}

\begin{problem}
若抛物线 \(y^2=2px\)（\(p>0\)）的焦点是椭圆 \(\dfrac{x^2}{3p}+\dfrac{y^2}{p}=1\) 的一个焦点，则 \(p=\)（\quad）

\choices
    {\(2\)}
    {\(3\)}
    {\(4\)}
    {\(8\)}
\end{problem}

\begin{answer}
D
\end{answer}

\begin{solution}
抛物线 \(y^2=2px\) 的焦点为 \(\left(\dfrac p2,0\right)\). 椭圆的半长轴、半短轴分别为 \(\sqrt{3p}\)、\(\sqrt p\)，所以其半焦距满足 \(c^2=3p-p=2p\)，焦点为 \(\left(\pm\sqrt{2p},0\right)\). 由两者有一个焦点重合且 \(p>0\)，得 \(\dfrac p2=\sqrt{2p}\)，所以 \(p=8\)，故选 D.
\end{solution}

\begin{problem}
曲线 \(y=2\sin x+\cos x\) 在点 \((\pi,-1)\) 处的切线方程为（\quad）

\choices
    {\(x-y-\pi-1=0\)}
    {\(2x-y-2\pi-1=0\)}
    {\(2x+y-2\pi+1=0\)}
    {\(x+y-\pi+1=0\)}
\end{problem}

\begin{answer}
C
\end{answer}

\begin{solution}
对 \(y=2\sin x+\cos x\) 求导，得 \(y'=2\cos x-\sin x\). 在 \(x=\pi\) 处，切线斜率为 \(y'(\pi)=-2\)，且切点为 \((\pi,-1)\). 因此切线方程为 \(y+1=-2(x-\pi)\)，整理得 \(2x+y-2\pi+1=0\)，故选 C.
\end{solution}

\begin{problem}
已知 \(\alpha\in\left(0,\dfrac{\pi}{2}\right)\)，\(2\sin2\alpha=\cos2\alpha+1\)，则 \(\sin\alpha=\)（\quad）

\choices
    {\(\dfrac15\)}
    {\(\dfrac{\sqrt5}{5}\)}
    {\(\dfrac{\sqrt3}{3}\)}
    {\(\dfrac{2\sqrt5}{5}\)}
\end{problem}

\begin{answer}
B
\end{answer}

\begin{solution}
由 \(\alpha\in\left(0,\dfrac{\pi}{2}\right)\)，有 \(\sin\alpha>0\) 且 \(\cos\alpha>0\). 将已知等式化为 \(4\sin\alpha\cos\alpha=\cos^2\alpha-\sin^2\alpha+1=2\cos^2\alpha\). 除以 \(2\cos\alpha\)，得 \(2\sin\alpha=\cos\alpha\). 结合 \(\sin^2\alpha+\cos^2\alpha=1\)，得 \(\sin\alpha=\dfrac1{\sqrt5}=\dfrac{\sqrt5}{5}\)，故选 B.
\end{solution}

\begin{problem}
设 \(F\) 为双曲线 \(C:\dfrac{x^2}{a^2}-\dfrac{y^2}{b^2}=1\)（\(a>0\)，\(b>0\)）的右焦点，\(O\) 为坐标原点，以 \(OF\) 为直径的圆与圆 \(x^2+y^2=a^2\) 交于 \(P\)，\(Q\) 两点. 若 \(|PQ|=|OF|\)，则 \(C\) 的离心率为（\quad）

\choices
    {\(\sqrt2\)}
    {\(\sqrt3\)}
    {\(2\)}
    {\(\sqrt5\)}
\end{problem}

\begin{answer}
A
\end{answer}

\begin{solution}
设双曲线的半焦距为 \(c\)，则 \(c^2=a^2+b^2\)，右焦点为 \(F=(c,0)\). 以 \(OF\) 为直径的圆的方程为 \(x^2+y^2=cx\). 该圆与 \(x^2+y^2=a^2\) 相交时，有 \(cx=a^2\)，所以交点的横坐标为 \(\dfrac{a^2}{c}\)，且 \(y^2=a^2-\dfrac{a^4}{c^2}=\dfrac{a^2b^2}{c^2}\). 因而弦长 \(PQ=\dfrac{2ab}{c}\). 由 \(PQ=OF=c\)，得 \(2ab=c^2=a^2+b^2\)，即 \((a-b)^2=0\)，所以 \(a=b\). 因此双曲线的离心率为 \(e=\dfrac ca=\dfrac{\sqrt{a^2+b^2}}a=\sqrt2\)，故选 A.
\end{solution}

\section{填空题}

\begin{problem}
若变量 \(x\)，\(y\) 满足约束条件 \(\begin{cases}
2x+3y-6\ge0,\\
x+y-3\le0,\\
y-2\le0,
\end{cases}\)，则 \(z=3x-y\) 的最大值是\fillinblank{}.
\end{problem}

\begin{answer}
\(9\)
\end{answer}

\begin{solution}
由 \(2x+3y\ge6\) 和 \(x+y\le3\)，可得 \(y\ge0\)；又已知 \(y\le2\). 由 \(x+y\le3\) 得 \(x\le3-y\)，所以 \(z=3x-y\le3(3-y)-y=9-4y\le9\). 点 \((3,0)\) 满足全部约束条件，且此时 \(z=9\)，故最大值为 \(9\).
\end{solution}

\begin{problem}
我国高铁发展迅速，技术先进. 经统计，在经停某站的高铁列车中，有 \(10\) 个车次的正点率为 \(0.97\)，有 \(20\) 个车次的正点率为 \(0.98\)，有 \(10\) 个车次的正点率为 \(0.99\)，则经停该站高铁列车所有车次的平均正点率的估计值为\fillinblank{}.
\end{problem}

\begin{answer}
\(0.98\)
\end{answer}

\begin{solution}
所有车次共有 \(10+20+10=40\) 个，因此平均正点率的估计值为 \(\dfrac{10\times0.97+20\times0.98+10\times0.99}{40}=\dfrac{39.2}{40}=0.98\).
\end{solution}

\begin{problem}
\(\triangle ABC\) 的内角 \(A\)，\(B\)，\(C\) 的对边分别为 \(a\)，\(b\)，\(c\). 已知 \(b\sin A+a\cos B=0\)，则 \(B=\fillinblank{}\).
\end{problem}

\begin{answer}
\(\dfrac{3\pi}{4}\)
\end{answer}

\begin{solution}
由正弦定理，设外接圆半径为 \(R\)，则 \(a=2R\sin A\)，\(b=2R\sin B\). 代入已知条件，得 \(2R\sin B\sin A+2R\sin A\cos B=2R\sin A(\sin B+\cos B)=0\). 由于三角形内角 \(A\in(0,\pi)\)，有 \(\sin A>0\)，所以 \(\sin B+\cos B=0\). 又 \(B\in(0,\pi)\)，故 \(B=\dfrac{3\pi}{4}\).
\end{solution}

\begin{problem}
中国有悠久的金石文化，印信是金石文化的代表之一. 印信的形状多为长方体，正方体或圆柱体，但南北朝时期的官员独孤信的印信形状是“半正多面体”（图1）. 半正多面体是由两种或两种以上的正多边形围成的多面体. 半正多面体体现了数学的对称美. 图2是一个棱数为 \(48\) 的半正多面体，它的所有顶点都在同一个正方体的表面上，且此正方体的棱长为 \(1\). 则该半正多面体共有\fillinblank{} 个面，其棱长为\fillinblank{}.

\examfiguregroup[float=inline]{img/2019/national_paper_2_liberal/q16_fig1.png}{%
    \bitmapinclude[max width=0.32\linewidth]{img/2019/national_paper_2_liberal/q16_fig1.png}\hfill
    \bitmapinclude[max width=0.40\linewidth]{img/2019/national_paper_2_liberal/q16_fig2.png}%
}
\end{problem}

\begin{answer}
\(26\)，\(\sqrt{2}-1\)
\end{answer}

\begin{solution}
由图 2 可知，每个顶点有 \(4\) 条棱相交. 设顶点数为 \(V\)，棱数为 \(E=48\)，则由每条棱有两个端点，得 \(4V=2E\)，所以 \(V=24\). 由欧拉公式 \(V-E+F=2\)，得面数 \(F=2-V+E=26\).

设正方体中心为原点. 由图形的对称性，半正多面体的顶点可表示为 \((\frac12,u,u)\) 的坐标排列并附以适当的正负号，其中 \(0<u<\frac12\). 图中同一正方体面的相邻两顶点可取为 \((\frac12,u,u)\) 和 \((\frac12,-u,u)\)，它们的距离为 \(2u\)；另一条相邻棱的两个端点可取为 \((\frac12,u,u)\) 和 \((u,\frac12,u)\)，其长度为 \(\sqrt2\left(\frac12-u\right)\). 半正多面体的所有棱长相等，因此 \(2u=\sqrt2\left(\dfrac12-u\right)\)，解得 \(u=\dfrac{\sqrt2-1}{2}\). 所以棱长为 \(2u=\sqrt2-1\).
\end{solution}

\section{解答题}

\begin{problem}
如图，长方体 \(ABCD-A_1B_1C_1D_1\) 的底面 \(ABCD\) 是正方形，点 \(E\) 在棱 \(AA_1\) 上，\(BE\perp EC_1\).

\begin{enumerate}
\item 证明：\(BE\perp\) 平面 \(EB_1C_1\)；
\item 若 \(AE=A_1E\)，\(AB=3\)，求四棱锥 \(E-BB_1C_1C\) 的体积.
\end{enumerate}

\bitmapfigure[max width=6.0cm]{img/2019/national_paper_2_liberal/q17_fig1.png}
\end{problem}

\begin{solution}
\begin{enumerate}
\item 设正方形的边长为 \(a\)，长方体的高为 \(h\)，建立直角坐标系，取 \(A=(0,0,0)\)，\(B=(a,0,0)\)，\(C=(a,a,0)\)，\(A_1=(0,0,h)\)，则 \(B_1=(a,0,h)\)，\(C_1=(a,a,h)\). 设 \(E=(0,0,t)\)，其中 \(0\le t\le h\)，
\(\overrightarrow{EB}=(a,0,-t)\)，\(\overrightarrow{EC_1}=(a,a,h-t)\). 由 \(BE\perp EC_1\)，得 \(\overrightarrow{EB}\cdot\overrightarrow{EC_1}=a^2-t(h-t)=0\)，即 \(t(h-t)=a^2\). 又 \(\overrightarrow{EB_1}=(a,0,h-t)\)，且 \(\overrightarrow{EB}\cdot\overrightarrow{EB_1}=a^2-t(h-t)=0\)，
所以 \(BE\perp EB_1\). 直线 \(EB_1\) 与 \(EC_1\) 是平面 \(EB_1C_1\) 内相交的两条直线，且 \(BE\) 同时垂直于它们，因此 \(BE\perp\) 平面 \(EB_1C_1\).

\item 由 \(AE=A_1E\) 得 \(t=h-t\)，即 \(t=h/2\). 结合 \(t(h-t)=a^2\)，有 \(\frac{h^2}{4}=a^2\). 由 \(AB=a=3\) 及 \(h>0\)，得 \(h=6\). 四棱锥的底面 \(BB_1C_1C\) 是矩形，其面积为 \(S_{BB_1C_1C}=BC\cdot BB_1=3\times6=18\). 点 \(E\) 到平面 \(BB_1C_1C\) 的距离等于 \(AB=3\)，所以四棱锥的体积为 \(V=\frac{1}{3}S_{BB_1C_1C}\cdot3=18\).
\end{enumerate}
\end{solution}

\begin{problem}
已知 \(\{a_n\}\) 是各项均为正数的等比数列，\(a_1=2\)，\(a_3=2a_2+16\).

\begin{enumerate}
\item 求 \(\{a_n\}\) 的通项公式；
\item 设 \(b_n=\log_2 a_n\)，求数列 \(\{b_n\}\) 的前 \(n\) 项和.
\end{enumerate}
\end{problem}

\begin{solution}
\begin{enumerate}
\item 设等比数列 \(\{a_n\}\) 的公比为 \(q\). 因为各项均为正数，所以 \(q>0\)，从而
\(a_2=2q\)，\(a_3=2q^2\).
代入已知条件，得
\(2q^2=4q+16\)，\(q^2-2q-8=0\)，\((q-4)(q+2)=0\).
由 \(q>0\)，得 \(q=4\)，因此
\(a_n=a_1q^{n-1}=2\cdot4^{n-1}\).

\item 由通项公式，
\[
b_n=\log_2a_n=\log_2\left(2\cdot4^{n-1}\right)=1+2(n-1)=2n-1
\]
所以数列 \(\{b_n\}\) 的前 \(n\) 项和为
\[
\sum_{k=1}^{n}b_k=\sum_{k=1}^{n}(2k-1)=2\cdot\frac{n(n+1)}{2}-n=n^2
\]
\end{enumerate}
\end{solution}

\begin{problem}
某行业主管部门为了解本行业中小企业的生产情况，随机调查了 \(100\) 个企业，得到这些企业第一季度相对于前一年第一季度产值增长率 \(y\) 的频数分布表.

\begin{center}
\begin{tabular}{c|ccccc}
\hline
\(y\) 的分组 & \([-0.20,0)\) & \([0,0.20)\) & \([0.20,0.40)\) & \([0.40,0.60)\) & \([0.60,0.80)\)\\
\hline
企业数 & 2 & 24 & 53 & 14 & 7\\
\hline
\end{tabular}
\end{center}

\begin{enumerate}
\item 分别估计这类企业中产值增长率不低于 \(40\%\) 的企业比例，产值负增长的企业比例；
\item 求这类企业产值增长率的平均数与标准差的估计值（同一组中的数据用该组区间的中点值为代表）. （精确到 \(0.01\)）
\end{enumerate}

附：\(\sqrt{74}\approx 8.602\).
\end{problem}

\begin{solution}
\begin{enumerate}
\item 产值增长率不低于 \(40\%\) 的企业有 \(14+7=21\) 个，所以这类企业的比例估计值为
\(\frac{14+7}{100}=0.21\). 产值负增长的企业有 \(2\) 个，所以这类企业的比例估计值为 \(\frac{2}{100}=0.02\).

\item 各组区间的中点依次为 \(-0.1\)，\(0.1\)，\(0.3\)，\(0.5\)，\(0.7\)，故产值增长率平均数的估计值为
\[
\bar y=\frac{2\times(-0.1)+24\times0.1+53\times0.3+14\times0.5+7\times0.7}{100}=0.30
\]
以这些中点值代替各组数据，标准差的估计值为
\[
\hat\sigma=\sqrt{\frac{2(0.4)^2+24(0.2)^2+53(0)^2+14(0.2)^2+7(0.4)^2}{100}}=\sqrt{0.0296}=\frac{\sqrt{74}}{50}\approx0.17
\]
因此，平均数与标准差的估计值分别为 \(0.30\) 和 \(0.17\).
\end{enumerate}
\end{solution}

\begin{problem}
已知 \(F_1\)，\(F_2\) 是椭圆 \(C:\dfrac{x^2}{a^2}+\dfrac{y^2}{b^2}=1\)（\(a>b>0\)）的两个焦点，\(P\) 为 \(C\) 上的点，\(O\) 为坐标原点.

\begin{enumerate}
\item 若 \(\triangle POF_2\) 为等边三角形，求 \(C\) 的离心率；
\item 如果存在点 \(P\)，使得 \(PF_1\perp PF_2\)，且 \(\triangle F_1PF_2\) 的面积等于 \(16\)，求 \(b\) 的值和 \(a\) 的取值范围.
\end{enumerate}
\end{problem}

\begin{solution}
\begin{enumerate}
\item 设椭圆的半焦距为 \(c\)，取 \(O=(0,0)\)，\(F_1=(-c,0)\)，\(F_2=(c,0)\). 因为 \(\triangle POF_2\) 为等边三角形，所以 \(PO=PF_2=OF_2=c\)，从而可取 \(P=\left(\frac{c}{2},\frac{\sqrt3c}{2}\right)\) 或 \(P=\left(\frac{c}{2},-\frac{\sqrt3c}{2}\right)\). 于是 \(PF_1=\sqrt3c\). 由椭圆的定义，\(PF_1+PF_2=2a\)，
故 \((\sqrt3+1)c=2a\). 因此离心率为
\[
e=\frac{c}{a}=\frac{2}{\sqrt3+1}=\sqrt3-1
\]

\item 设 \(u=PF_1\)，\(v=PF_2\). 由椭圆的定义及勾股定理，得
\(u+v=2a\)，\(u^2+v^2=F_1F_2^2=4c^2\).
又因为 \(PF_1\perp PF_2\)，且三角形 \(F_1PF_2\) 的面积为 \(16\)，所以
\(\frac12uv=16\)，\(uv=32\).
于是
\[
4a^2=(u+v)^2=u^2+v^2+2uv=4c^2+64
\]
即 \(a^2-c^2=16\). 由 \(b^2=a^2-c^2\)，得 \(b=4\).

为使这样的点 \(P\) 存在，\(u\)，\(v\) 必须为正数，因此方程
\(t^2-2at+32=0\)
必须有正实根，故
\(4a^2-128\ge0\)，\(a\ge4\sqrt2\).
反之，当 \(a\ge4\sqrt2\) 时，取 \(u=a+\sqrt{a^2-32}\)，\(v=a-\sqrt{a^2-32}\)，则 \(u\)，\(v>0\)，且
\(u+v=2a\)，\(uv=32\)，\(u^2+v^2=4(a^2-16)=4c^2\).
因此以 \(F_1F_2=2c\) 为一边、另外两边为 \(u\)，\(v\) 的三角形存在，并且在 \(P\) 处为直角三角形，其顶点 \(P\) 满足椭圆定义. 所以 \(a\) 的取值范围为
\(a\in[4\sqrt2,+\infty)\).
\end{enumerate}
\end{solution}

\begin{problem}
已知函数 \(f(x)=(x-1)\ln x-x-1\). 证明：

\begin{enumerate}
\item \(f(x)\) 存在唯一的极值点；
\item \(f(x)=0\) 有且仅有两个实根，且两个实根互为倒数.
\end{enumerate}
\end{problem}

\begin{solution}
\begin{enumerate}
\item 函数的定义域为 \((0,+\infty)\)，且
\(f'(x)=\ln x+\frac{x-1}{x}-1=\ln x-\frac1x\)，\(f''(x)=\frac1x+\frac1{x^2}=\frac{x+1}{x^2}>0\).
所以 \(f'(x)\) 在 \((0,+\infty)\) 上严格递增. 又
\(f'(1)=-1<0\)，\(f'(\e)=1-\frac1{\e}>0\)，
由零点存在性可知，存在唯一的 \(x_0\in(1,\e)\)，使得 \(f'(x_0)=0\). 因而当 \(0<x<x_0\) 时 \(f'(x)<0\)，当 \(x>x_0\) 时 \(f'(x)>0\)，所以 \(f(x)\) 在 \(x_0\) 处取得唯一的极小值，故存在唯一的极值点.

\item 当 \(x\to0^+\) 时，\((x-1)\ln x\to+\infty\)，故 \(f(x)\to+\infty\). 又当 \(x\to+\infty\) 时，\(f(x)\to+\infty\)，且 \(f(1)=-2<0\). 结合上一问的单调性，\(f(x)\) 在 \((0,x_0)\) 上严格递减，在 \((x_0,+\infty)\) 上严格递增，因此在 \((0,1)\) 内有且仅有一个零点，在 \((x_0,+\infty)\) 内有且仅有一个零点，故 \(f(x)=0\) 有且仅有两个实根.

对任意 \(x>0\)，有
\[
f\left(\frac1x\right)=\left(\frac1x-1\right)\ln\frac1x-\frac1x-1=\frac{f(x)}{x}
\]
因此若 \(x\) 是方程 \(f(x)=0\) 的根，则 \(1/x\) 也是该方程的根. 由于方程恰有两个根，且一个根在 \((0,1)\) 内、另一个根大于 \(1\)，这两个根互为倒数.
\end{enumerate}
\end{solution}

\section{选考题：解答题}

\begin{problem}
在极坐标系中，\(O\) 为极点，点 \(M(\rho_0,\theta_0)\)（\(\rho_0>0\)）在曲线 \(C:\rho=4\sin\theta\) 上，直线 \(l\) 过点 \(A(4,0)\) 且与 \(OM\) 垂直，垂足为 \(P\).

\begin{enumerate}
\item 当 \(\theta_0=\dfrac{\pi}{3}\) 时，求 \(\rho_0\) 及 \(l\) 的极坐标方程；
\item 当 \(M\) 在 \(C\) 上运动且 \(P\) 在线段 \(OM\) 上时，求 \(P\) 点轨迹的极坐标方程.
\end{enumerate}
\end{problem}

\begin{solution}
\begin{enumerate}
\item 由点 \(M(\rho_0,\theta_0)\) 在曲线 \(C\) 上，得
\(\rho_0=4\sin\theta_0=4\sin\frac{\pi}{3}=2\sqrt3\).
设直线 \(l\) 上任意一点的极坐标为 \((\rho,\theta)\). 因为 \(l\perp OM\)，直线 \(OM\) 的方向向量为 \((\cos\theta_0,\sin\theta_0)\)，所以
\(x\cos\theta_0+y\sin\theta_0=4\cos\theta_0\).
代入 \(x=\rho\cos\theta\)，\(y=\rho\sin\theta\)，得直线 \(l\) 的极坐标方程
\(\rho\cos(\theta-\theta_0)=4\cos\theta_0\).
当 \(\theta_0=\dfrac{\pi}{3}\) 时，
\(\rho\cos\left(\theta-\frac{\pi}{3}\right)=2\).

\item 由于 \(P\) 在线段 \(OM\) 上，除 \(P=O\) 外，点 \(P\) 与点 \(M\) 的极角相同，记为 \(\theta\). 设 \(P=(\rho,\theta)\)，由 \(AP\perp OP\) 得
\(\bigl(4-\rho\cos\theta,-\rho\sin\theta\bigr)\cdot(\cos\theta,\sin\theta)=0\)，
从而 \(\rho=4\cos\theta\).

又因为 \(M\) 在曲线 \(C\) 上且 \(\rho_0>0\)，有 \(\rho_0=4\sin\theta>0\). 由于 \(P\) 在线段 \(OM\) 上，还需满足
\(0\le\rho\le\rho_0\)，\(0\le4\cos\theta\le4\sin\theta\).
结合 \(0<\theta<\pi\)，可得 \(\dfrac{\pi}{4}\le\theta\le\dfrac{\pi}{2}\). 因此 \(P\) 点的轨迹为
\(\rho=4\cos\theta\)，\(\frac{\pi}{4}\le\theta\le\frac{\pi}{2}\).
\end{enumerate}
\end{solution}

\begin{problem}
已知 \(f(x)=|x-a|x+|x-2|(x-a)\).

\begin{enumerate}
\item 当 \(a=1\) 时，求不等式 \(f(x)<0\) 的解集；
\item 当 \(x\in(-\infty,1)\) 时，\(f(x)<0\)，求 \(a\) 的取值范围.
\end{enumerate}
\end{problem}

\begin{solution}
\begin{enumerate}
\item 当 \(a=1\) 时，按 \(x<1\)，\(1\le x<2\)，\(x\ge2\) 分段讨论.

当 \(x<1\) 时，
\[
f(x)=(1-x)x+(2-x)(x-1)=-2(x-1)^2<0
\]
当 \(1\le x<2\) 时，
\[
f(x)=x(x-1)+(2-x)(x-1)=2(x-1)\ge0
\]
当 \(x\ge2\) 时，
\[
f(x)=x(x-1)+(x-2)(x-1)=2(x-1)^2>0
\]
所以不等式 \(f(x)<0\) 的解集为 \(\left(-\infty,1\right)\).

\item 要使任意 \(x<1\) 时都有 \(f(x)<0\)，先证明 \(a\ge1\). 若 \(a<1\)，取 \(x=\dfrac{a+1}{2}\)，则 \(a<x<1\)，此时 \(|x-a|=x-a\)，且 \(|x-2|=2-x\)，从而
\[
f(x)=(x-a)x+(2-x)(x-a)=2(x-a)>0
\]
与题设矛盾.

当 \(a\ge1\) 时，对任意 \(x<1\)，都有 \(x<a\)，于是
\[
f(x)=(a-x)x+(2-x)(x-a)=-2(x-1)(x-a)=2(1-x)(x-a)<0
\]
因此 \(a\) 的取值范围为 \([1,+\infty)\).
\end{enumerate}
\end{solution}
